\scanpage{34}
\mspage{6 - 36 -}
n\textsuperscript{o} 392

g) Soient $K$ un corps, $A$ l'anneau de polynômes\unsure{.} $K[X,Y,Z]$, $\underline m$ l'idéal maximal $AX+AY+AZ$, $\underline q$ l'idéal de $A$ engendré par $X$, $Y^2$, $YZ$ et $Z^2$. On considère le $A$-module $M=A/\underline q$ et son sous-module $N=\underline m/\underline q$. Montrer que l'on a $\underline d_o(M)=\underline q$, $\underline d_1(M)=A$, $\underline d_o(N)=\underline m^2$, $\underline d_1(N)=\underline m$.

11) Soient $A$ un anneau de Krull, $M$, $N$ deux réseaux de type fini dans un espace vectoriel $V$ de dimension finie $n$ sur le corps des fractions de $A$. Pour tout $c\ne 0$ dans $A$ tel que $cN\subset M$, on considère les idéaux déterminantiels $\underline d_k(M/cN)$, (exerc. 10), et on pose :
\[
 d_k(N,M)=\operatorname{div}(\underline d_k(M/cN))-(n-k)\operatorname{div}(c).
\]

a) Montrer que $d_k(N,M)$ ne dépend pas du choix de $c$ tel que $cN\subset M$ ; on dit que $d_k(N,M)$ est le \uline{diviseur déterminantiel} d'indice $k$ de $N$ par rapport à $M$.

b) On a $d_k(N,M)\geq d_{k+1}(N,M)$, et $d_n(N,M)=0$. Montrer que si l'on pose $e_k(N,M)=d_{n-k}(N,M)-d_{n-k+1}(N,M)$, on a $e_k(N,M)\leq e_{k+1}(N,M)$ pour $1\leq k\leq n$ ; on dit que les diviseurs $e_k(N,M)$ sont les \uline{facteurs invariants} de $N$ par rapport à $M$, (se ramener au cas o\rlap{\handwritten{\`{\phantom{u}}}}u $A$ est un anneau principal).

c) Montrer que l'on a $d_o(N,M)=\chi(M,N)$, (n\textsuperscript{o}\quad).

d) Lorsque $M$ est un réseau dans $V$ mais que $N$ est un réseau dans un sous-espace $W$ de $V$ de dimension $q<n$, on appelle \uline{facteurs invariants} de $N$ par rapport à $M$ les facteurs invariants de $N$ par rapport à $M\cap W$. Montrer comment s'étendent les exerc. 8 à 10 et 14 à 16 d'\uline{Alg}., chap. VII, \S~4, (en supposant au besoin dans certains cas que $A$ est noethérien et intégralement clos).

12) Soient $A$ un anneau noethérien intégralement clos, $M$, $N$ deux $A$-modules sans torsion de type fini, de même rang $r$, tels que $N\subset M$ ; soit $j:N\longrightarrow M$ l'injection canonique.

a) Montrer que, pour tout $K$, $\Lambda^k j:\Lambda^k N\longrightarrow\Lambda^k M$ est pseudo-injectif, que $\operatorname{Coker}(\Lambda^k j)$ est un $A$-module de torsion et que l'on a :
\[
 \chi(\operatorname{Coker}(\Lambda^k j))=\binom{r-1}{k-1}\chi(\operatorname{Coker}j).
\]

b) En utilisant a), démontrer que si $M$ est un $A$-module sans torsion de type fini, le sous-module de torsion de $\Lambda^k M$ est pseudo-nul pour tout $k$, et l'on a $c(\Lambda^k M)=\binom{r-1}{k-1}c(M)$, o\rlap{\handwritten{\`{\phantom{u}}}}u $M$ est de rang $r$ ; en particulier, on a $c(\Lambda^r M)=c(M)$.
% archive sheet number: 33

